%!TEX root = ./main.tex

\section[Wrap-up]{Wrap-Up}

% SECTION TIMING: 43:10--50:00 (3 frames): 150 + 160 + 100 = 410 s.
% If behind, show only the first two points of the open-problems frame (90 s).

% Teaching note (150 s): retrieval, not new material -- before each click, ask the
% class to fill in the next row. Rows: today's pipeline reacts one step late and
% sends everything ("Where it struggles"); Version 1 needs the same model on both
% ends and a wrong frame looks fine ("The price of rebuilding"); Version 2 is built
% for the tasks it was trained for -- a new kind of task may need retraining, and
% general multi-task models exist (e.g., U-DeepSC, IEEE Trans. Commun. 2024) but send
% more bits and need more compute -- and there is no video to replay; Version 3 works
% with today's codecs and networks but still sends a picture and needs AI on both
% ends ("Version 3 in practice"). The takeaway answers "which one wins?": none --
% choose by what the receiver needs.
\begin{frame}{Four ways to send video}
  \vspace{0.8em}
  \centering
  {\footnotesize\renewcommand{\arraystretch}{1.45}\setlength{\tabcolsep}{3pt}%
  \begin{tabular}{@{}>{\bfseries}l l l l l@{}}
    \toprule
    & \textbf{Sends} & \textbf{Receiver gets} & \textbf{Strength} & \textbf{Weak spot} \\
    \midrule
    \textcolor{ranblue}{Today}
      & every bit & the exact video & works for any app
      & \makecell[lt]{reacts late;\\sends everything} \\
    \uncover<2->{\textcolor{ranpurple}{Version 1}}
      & \uncover<2->{meaning} & \uncover<2->{a close copy}
      & \uncover<2->{degrades smoothly}
      & \uncover<2->{\makecell[lt]{needs same model;\\errors look fine}} \\
    \uncover<3->{\textcolor{coreorange}{Version 2}}
      & \uncover<3->{task clues} & \uncover<3->{the answer}
      & \uncover<3->{far fewer bits}
      & \uncover<3->{\makecell[lt]{built for its tasks;\\no video to replay}} \\
    \uncover<4->{\textcolor{softgreen}{Version 3}}
      & \uncover<4->{\makecell[lt]{bits, spent\\by meaning}}
      & \uncover<4->{\makecell[lt]{sharp where\\it matters}}
      & \uncover<4->{\makecell[lt]{works on today's\\networks}}
      & \uncover<4->{\makecell[lt]{still sends a picture;\\AI on both ends}} \\
    \bottomrule
  \end{tabular}\par}

  \vspace{1.2em}
  \takeawayat{5-}{No single winner:\\
    choose by what the receiver needs.}
\end{frame}

% Teaching note (160 s): status check, one click per point. Semantic communication
% is research-stage. (1) Today's radios and protocols assume bits (fixed signal
% alphabets, retransmissions, layered stacks). (2) Who trains, ships, and updates the
% shared model on every device? (3) Trust: no checksum for meaning; crafted inputs can
% fool the AI models; transmitted features can still leak private content.
% (4) Standards, checked to about June 2026: 3GPP's AI/ML air-interface work (Rel-18
% study TR 38.843; Rel-19 normative one-sided models) targets CSI feedback/prediction,
% beam management, and positioning; two-sided CSI compression is deferred to Rel-20;
% there is no 3GPP study or work item on semantic communication, and the 6G use-case
% report TR 22.870 mentions it only as motivation inside use cases. 3GPP's other AI
% work (core-network analytics such as NWDAF, AI/ML for NG-RAN) is not about content
% transmission either. Recheck before each delivery of this lecture.
% Version 3 sidesteps the first and last points by design (standard codecs, standard
% networks), which is why it is the closest to deployment today.
% MPEG (the video-codec standards body) is closer to Version 2: Video Coding for
% Machines, VCM (ISO/IEC 23888-2), codes video meant for machine analysis -- its
% decoder still outputs video -- and reached Draft International Standard at MPEG 153
% (January 2026), with the final stage planned for late 2026; Feature Coding for
% Machines, FCM (ISO/IEC 23888-4), codes the intermediate features of a split neural
% network so the cloud "only needs to run the head", and reached Committee Draft at
% MPEG 154 (April-May 2026). These are source-coding standards (no joint design with
% the radio channel). Sources: MPEG press releases for meetings 153-155; stage data
% from the ISS mirror of the ISO record (iso.org itself could not be fetched).
% Optional systems examples if students ask "who builds this?": MagicStream (ACM
% SenSys 2024, semantic 3D telepresence), KDC (USENIX NSDI 2026, knowledge-driven
% generative communication), LITE (ACM SIGCOMM 2026, loss-resilient telepresence with
% multi-modal semantics).
\begin{frame}{Why it is not in your phone yet}
  \vspace{1.2em}
  \begin{enumerate}[<+->]\itemsep12pt
    \item \textbf{Must fit today's radios and networks}
    \item \textbf{Both sides must keep the same AI model}
    \item \textbf{Trust:} no checksum for meaning; AI can be fooled
    \item \textbf{Standards are just starting}\\
      {\small\textcolor{softgray}{no cellular standard yet; MPEG is defining video and
        feature coding \textit{for machines} (2026)}}
  \end{enumerate}

  \vspace{1.6em}
  \takeawayat{5-}{Promising for weak signals and clear goals;\\
    not a replacement for today's networks yet.}
  \source{https://portal.3gpp.org/desktopmodules/Specifications/SpecificationDetails.aspx?specificationId=3983}{3GPP
    TR 38.843 (Rel-18); MPEG press releases, meetings 153--155 (2026)}
\end{frame}

% Teaching note (100 s): four sentences to remember, one click each, then close by
% re-asking the opening question: A needs (a good copy of) the video, B needs only
% the answer. (The pair discussion was removed to make room for Version 3.)
\begin{frame}{Takeaways}
  \vspace{1.2em}
  \begin{enumerate}[<+->]\itemsep10pt
    \item \textbf{\textcolor{ranblue}{Today:}} copy every bit of every frame
    \item \textbf{\textcolor{ranpurple}{Version 1:}} rebuild the video;
      degrades smoothly
    \item \textbf{\textcolor{coreorange}{Version 2:}} no video rebuilt ---
      send only what the task needs
    \item \textbf{\textcolor{softgreen}{Version 3:}} today's pipe,
      bits spent by meaning
  \end{enumerate}

  \vspace{2.0em}
  \askat{5-}{Back to the start: do A and B need the same data?}
\end{frame}
